Considera la matriu $M=\begin{pmatrix}1&a&1\\a&1&1\\1&1&a\end{pmatrix}$, en què $a$ és un nombre real.

\begin{apartats}

\apartat[1,25]{0,75}
Calcula $|M|$ i comprova que és igual a $-(a-1)^2(a+2)$.

\begin{solucio}
Per Sarrus: $|M|=a+a+a-1-1-a^3=-a^3+3a-2$.\\
D'altra banda, $-(a-1)^2(a+2)=-(a^2-2a+1)(a+2)=-(a^3-3a+2)=-a^3+3a-2$: coincideixen.
\end{solucio}

\begin{tria}{rang-parametre}
\itemtria{estudia-rang}{1}{1,25}
Estudia el rang de $M$ segons els valors de $a$.

\begin{solucio}
\begin{itemize}
\item Si $a\ne1$ i $a\ne-2$, $|M|\ne0$: $\operatorname{rang}M=3$.
\item Si $a=1$, les tres files són $(1\ \ 1\ \ 1)$: $\operatorname{rang}M=1$.
\item Si $a=-2$, $M=\begin{pmatrix}1&-2&1\\-2&1&1\\1&1&-2\end{pmatrix}$, $|M|=0$ i
$\begin{vmatrix}1&-2\\-2&1\end{vmatrix}=-3\ne0$: $\operatorname{rang}M=2$.
\end{itemize}
\end{solucio}

\itemtria{rang-1}{1}{1,25}
Troba els valors de $p$ i $q$ perquè la matriu $\begin{pmatrix}1&p&3\\2&4&q\end{pmatrix}$ tingui rang 1.

\begin{solucio}
La matriu no és nul·la, i té rang 1 si i només si tots els seus menors d'ordre 2 són nuls:
$\begin{vmatrix}1&p\\2&4\end{vmatrix}=4-2p=0$ dona $p=2$, i $\begin{vmatrix}1&3\\2&q\end{vmatrix}=q-6=0$
dona $q=6$. Amb $p=2$ i $q=6$, el tercer menor, $\begin{vmatrix}2&3\\4&6\end{vmatrix}$, també és nul: la
segona fila és el doble de la primera.
\end{solucio}
\end{tria}

\begin{nomesllarg}
\apartat{0,75}
Per a $a=-2$, troba una combinació lineal de les files de $M$ que doni la fila nul·la.

\begin{solucio}
Amb $a=-2$, les files són $F_1=(1\ \ {-2}\ \ 1)$, $F_2=(-2\ \ 1\ \ 1)$ i $F_3=(1\ \ 1\ \ {-2})$.
$F_1+F_2+F_3=(0\ \ 0\ \ 0)$: la tercera fila és $F_3=-F_1-F_2$, i per això el rang no arriba a 3.
\end{solucio}
\end{nomesllarg}

\end{apartats}
